\documentclass[11pt]{article}

% ============================================================
% Plan de clase — Semana 4 (uso exclusivo del docente)
% Ingeniería de Requerimientos (ISR-401) · UTEQ · SPA 2026-2027
% Compilación: pdflatex -> bibtex -> pdflatex -> pdflatex
% ============================================================

\usepackage[utf8]{inputenc}
\usepackage[T1]{fontenc}
\usepackage{mathptmx}
\usepackage[scaled]{helvet}
\usepackage{courier}
\renewcommand{\familydefault}{\sfdefault}
\usepackage[a4paper,top=2.6cm,bottom=2.3cm,left=2cm,right=2cm,headheight=34pt]{geometry}
\usepackage{graphicx}
\usepackage{array}
\usepackage{tabularx}
\usepackage{longtable}
\usepackage{colortbl}
\usepackage{xcolor}
\usepackage{enumitem}
\usepackage{fancyhdr}
\usepackage{lastpage}
\usepackage[most]{tcolorbox}
\usepackage{amssymb}
\usepackage[numbers,sort&compress]{natbib}
\usepackage{url}
\usepackage{needspace}
\usepackage{caption}
\usepackage[hidelinks]{hyperref}

\renewcommand{\tablename}{Tabla}
\renewcommand{\refname}{Referencias}

\definecolor{utegreen}{HTML}{0B7A3B}
\definecolor{utegreendk}{HTML}{075C2C}
\definecolor{midblue}{HTML}{1F5C8B}
\definecolor{gold}{HTML}{D4A017}
\definecolor{rowgray}{HTML}{EFEFEF}
\definecolor{boxgray}{HTML}{F4F6F4}
\definecolor{lightred}{HTML}{FBEAEA}
\definecolor{lightgold}{HTML}{FBF3DC}
\definecolor{lightblue}{HTML}{EAF1F7}

\newcolumntype{L}[1]{>{\raggedright\arraybackslash}p{#1}}
\newcolumntype{C}[1]{>{\centering\arraybackslash}p{#1}}
\newcolumntype{Y}{>{\raggedright\arraybackslash}X}

\pagestyle{fancy}
\fancyhf{}
\lhead{\footnotesize UTEQ \,$\cdot$\, Ingeniería de Requerimientos (ISR-401)}
\rhead{\footnotesize Plan de clase de la semana 4 \,$\cdot$\, Uso del docente}
\lfoot{\footnotesize Prof. Gleiston C. Guerrero-Ulloa, PhD}
\rfoot{\footnotesize Página \thepage\ de \pageref{LastPage}}
\renewcommand{\headrulewidth}{0.6pt}
\renewcommand{\footrulewidth}{0.4pt}

\setlength{\parindent}{0pt}
\setlength{\parskip}{5pt}
\setlength{\emergencystretch}{3em}
\renewcommand{\arraystretch}{1.25}

\newcommand{\thdr}[1]{\cellcolor{utegreen}\textcolor{white}{\textbf{#1}}}
\newcommand{\seccion}[1]{\section{\textcolor{utegreendk}{#1}}}
\renewcommand{\thesection}{\arabic{section}.}
\newtcolorbox{cajaclave}[1][]{colback=boxgray,colframe=utegreen,boxrule=0.8pt,arc=2mm,fonttitle=\bfseries,breakable,#1}
\newtcolorbox{cajaaviso}[1][]{colback=lightgold,colframe=gold!80!black,boxrule=0.8pt,arc=2mm,fonttitle=\bfseries,breakable,#1}
\newcommand{\caja}{$\square$}

\begin{document}

\noindent\includegraphics[width=\linewidth]{logoband.png}

\begin{center}
	{\LARGE\textbf{\textcolor{utegreendk}{PLAN DE CLASE DE LA SEMANA 4}}}\\[4pt]
	{\large Unidad 2, Tema 1: Elicitación de requisitos}\\[4pt]
	{\small Documento interno: no se publica a los estudiantes}
\end{center}
\noindent\textcolor{gold}{\rule{\linewidth}{1.5pt}}

% ============================================================
\seccion{Datos de la semana según el sílabo}
% ============================================================

\noindent\begin{tabularx}{\linewidth}{|L{0.30\linewidth}|Y|}
	\hline
	\thdr{Campo} & \thdr{Detalle}\\
	\hline
	Semana y fechas & 4: lunes 16 a domingo 22 de noviembre de 2026 (sin feriados)\\
	\hline
	\rowcolor{rowgray} Unidad & 2. Elicitación y Análisis de Requisitos\\
	\hline
	Resultado de aprendizaje & Aplica técnicas de elicitación y análisis de requisitos para identificar, clasificar y modelar requisitos.\\
	\hline
	\rowcolor{rowgray} Contenidos (Tema 1) & Fuentes de requisitos: stakeholders, documentos, sistemas existentes y estándares. Técnicas: entrevista, taller (workshop), focus group, observación y etnografía. Técnicas: cuestionario, lectura basada en perspectivas, design thinking. Técnicas de apoyo: brainstorming, prototipado, KJ, mapas mentales y checklists.\\
	\hline
	Estrategias metodológicas & Aprendizaje basado en problemas, aula invertida y design thinking.\\
	\hline
	\rowcolor{rowgray} Habilidades blandas & Comunicación efectiva, creatividad y empatía.\\
	\hline
	TAC & LMS (SGA), herramientas de ideación y prototipado, plataformas de comunicación y gestión grupal.\\
	\hline
	\rowcolor{rowgray} Recursos didácticos & Cuestionario en línea, libros de texto, material audiovisual y presentaciones.\\
	\hline
	Evaluación & GA-TL-02 (formativa, taller) y TA-PFC-E1 (sumativa, 1.75 puntos del corte 1). Ambas cierran el sábado 21 de noviembre de 2026 a las 23:59.\\
	\hline
	\rowcolor{rowgray} Horas & 3 en contacto con el docente, 2 práctico-experimentales y 4 autónomas.\\
	\hline
\end{tabularx}

% ============================================================
\seccion{Material de la semana}
% ============================================================

{\small
	\begin{longtable}{|L{0.40\linewidth}|L{0.14\linewidth}|L{0.38\linewidth}|}
		\caption{Paquete de la semana 4 y su destino.}\\
		\hline
		\thdr{Archivo} & \thdr{Para} & \thdr{Uso}\\
		\hline
		\endfirsthead
		\hline
		\thdr{Archivo} & \thdr{Para} & \thdr{Uso}\\
		\hline
		\endhead
		\texttt{Semana04\_\allowbreak{}Elicitacion\_\allowbreak{}ISR401.pdf} (y .tex) & Estudiantes & Lectura de aula invertida (16 páginas); se publica en el SGA el lunes 16 antes de la primera sesión.\\
		\hline
		\rowcolor{rowgray} \texttt{ISR401\_\allowbreak{}Semana04.pptx} & Docente y estudiantes & Diapositivas de las dos sesiones (20 diapositivas, con notas del orador). Reemplaza a la versión anterior, que no correspondía al sílabo.\\
		\hline
		\texttt{Plantilla\_\allowbreak{}GA-TL-02/FORM/TL-02/} & Estudiantes & Plantilla \LaTeX{} del taller (\texttt{main.tex}, \texttt{referencias.bib}, \texttt{evidencias/MANIFIESTO.md}, instrucciones de compilación) y su PDF; la página del mapa de empatía sirve para imprimir.\\
		\hline
		\rowcolor{rowgray} \texttt{Guia\_\allowbreak{}Rubrica\_\allowbreak{}GA-TL-02\_\allowbreak{}Entrevista\_\allowbreak{}DesignThinking.pdf} & Estudiantes & Guía y lista de cotejo del taller (ya existente, sin cambios).\\
		\hline
		\texttt{Guia\_\allowbreak{}Rubrica\_\allowbreak{}TA-PFC-E1\_\allowbreak{}Planificacion\_\allowbreak{}Elicitacion.pdf} & Estudiantes & Guía y rúbrica de la entrega 1 del PFC (ya publicada, sin cambios).\\
		\hline
		\rowcolor{rowgray} \texttt{Consentimiento\_\allowbreak{}Informado\_\allowbreak{}PFC\_\allowbreak{}ISR401.pdf} & Estudiantes & Modelo de consentimiento para el ensayo y para el trabajo de campo del PFC (ya existente).\\
		\hline
		\texttt{referencias\_\allowbreak{}semana04.bib} & Docente & Referencias verificadas de la semana, en BibTeX.\\
		\hline
		\rowcolor{rowgray} Este plan & Docente & Secuencia, cuestionario con clave y observaciones de alineación.\\
		\hline
	\end{longtable}
}

\begin{cajaclave}[title={Lista de verificación antes del lunes 16}]
	\begin{itemize}[leftmargin=18pt,itemsep=2pt,label=\caja]
		\item Publicar en el SGA la lectura y el cuestionario de aula invertida (sección 5), con cierre antes de la primera sesión.
		\item Confirmar que las guías de GA-TL-02 y TA-PFC-E1 están visibles en el SGA y publicar la plantilla de GA-TL-02.
		\item Imprimir una hoja del mapa de empatía (página 3 de la plantilla) por equipo y llevar notas adhesivas para el método KJ.
		\item Llevar copias del modelo de consentimiento informado para el ensayo en juego de roles.
		\item Verificar qué equipos ya tienen la carta de aceptación de su organización (TA-PFC-E1).
	\end{itemize}
\end{cajaclave}

% ============================================================
\seccion{Sesión de contacto con el docente (3 horas)}
% ============================================================

Estrategia: aula invertida y aprendizaje basado en problemas. El problema conductor es el sistema del PFC de cada equipo: cada concepto se aplica inmediatamente a la organización real del equipo.

{\small
	\begin{longtable}{|C{0.08\linewidth}|L{0.30\linewidth}|L{0.38\linewidth}|C{0.13\linewidth}|}
		\caption{Secuencia de la sesión de contacto.}\\
		\hline
		\thdr{Min.} & \thdr{Actividad} & \thdr{Qué hacen docente y estudiantes} & \thdr{Diapos.}\\
		\hline
		\endfirsthead
		\hline
		\thdr{Min.} & \thdr{Actividad} & \thdr{Qué hacen docente y estudiantes} & \thdr{Diapos.}\\
		\hline
		\endhead
		0--15 & Inicio y cuestionario & Agenda y resultado de aprendizaje; revisión de las preguntas del cuestionario con más errores. & 1--3\\
		\hline
		\rowcolor{rowgray} 15--45 & Dificultad y fuentes & Conocimiento tácito con el ejemplo del cajero; cada equipo nombra una fuente de cada tipo para su PFC y tres roles de partes interesadas. & 4--6\\
		\hline
		45--55 & Mapa de técnicas & Clasificación IREB; los estudiantes ubican cada técnica del sílabo. & 7\\
		\hline
		\rowcolor{rowgray} 55--100 & Técnicas de recolección & Entrevista y tipos de preguntas; ejercicio de 5 minutos: cada equipo escribe una pregunta defectuosa sobre su PFC y el equipo vecino la corrige. Cuestionario, taller, grupo focal, observación, etnografía y lectura basada en perspectivas. & 8--12\\
		\hline
		100--110 & Receso & --- & ---\\
		\hline
		\rowcolor{rowgray} 110--150 & Design thinking y apoyo & Fases y mapa de empatía; ejemplo bueno y malo de «¿Cómo podríamos\ldots?». Dinámica KJ de 10 minutos con notas adhesivas sobre problemas típicos de una cooperativa. & 13--14\\
		\hline
		150--170 & Selección, IA y ética & Tabla de selección; cada equipo elige la técnica para tres de sus partes interesadas. Demostración de \texttt{Get-FileHash} y del manifiesto. & 15--16\\
		\hline
		\rowcolor{rowgray} 170--180 & Cierre & Requisitos de TA-PFC-E1 y anuncio del taller; dudas. & 18--19\\
		\hline
	\end{longtable}
}

% ============================================================
\seccion{Sesión práctico-experimental (2 horas): taller GA-TL-02}
% ============================================================

Estrategia: design thinking. Los equipos trabajan en su repositorio con la plantilla. Para que el ensayo se parezca a una entrevista real, la parte interesada la representa un integrante de otro equipo, que no conoce el guion.

{\small
	\begin{longtable}{|C{0.10\linewidth}|L{0.25\linewidth}|L{0.55\linewidth}|}
		\caption{Secuencia del taller.}\\
		\hline
		\thdr{Min.} & \thdr{Paso} & \thdr{Indicaciones}\\
		\hline
		\endfirsthead
		\hline
		\thdr{Min.} & \thdr{Paso} & \thdr{Indicaciones}\\
		\hline
		\endhead
		0--10 & Organización & Repaso de la diapositiva 17; emparejar equipos (A entrevista a un integrante de B y viceversa).\\
		\hline
		\rowcolor{rowgray} 10--40 & 1. Guion & 12 a 15 preguntas con apertura, bloques y cierre, marcando el tipo de cada una. El docente recorre los equipos y señala preguntas inductivas o dobles sin corregirlas.\\
		\hline
		40--80 & 2. Ensayo & Dos rondas de 20 minutos (cada equipo entrevista y es entrevistado). Lectura y firma del consentimiento; grabación solo de audio; un integrante pregunta y otro registra.\\
		\hline
		\rowcolor{rowgray} 80--110 & 3. Mapa y problema & Mapa de empatía en papel con lo dicho en el ensayo; declaración del problema y pregunta «¿Cómo podríamos\ldots?»; foto del mapa a \texttt{figuras/}.\\
		\hline
		110--120 & 4. Ajustes y commit & Tabla de preguntas que fallaron y su corrección; primer commit del día en \texttt{FORM/TL-02/}.\\
		\hline
	\end{longtable}
}

\begin{cajaaviso}[title={Puntos a observar durante el taller (lista de cotejo de GA-TL-02)}]
	\begin{enumerate}[leftmargin=18pt,itemsep=2pt]
		\item Guion con estructura y preguntas abiertas, neutrales y simples.
		\item Ensayo con consentimiento y registrado en el manifiesto (no audios en el repositorio).
		\item Mapa de empatía coherente con lo dicho en el ensayo, no con lo que el equipo supone.
		\item Definición del problema centrada en el usuario, no en la solución.
		\item Preguntas defectuosas identificadas y corregidas.
	\end{enumerate}
	La lista de cotejo se devuelve en el SGA antes de la sesión siguiente, para que los indicadores EP o NL se corrijan antes de GA-PE-02 y TA-PFC-02.
\end{cajaaviso}

% ============================================================
\seccion{Cuestionario de aula invertida (formativo, SGA)}
% ============================================================

Diez preguntas de opción múltiple sobre la lectura, con una sola respuesta correcta (marcada con \textbf{*}). Configuración sugerida: un intento, 20 minutos, cierre antes de la primera sesión de la semana, sin calificación.

\begin{enumerate}[leftmargin=20pt,itemsep=6pt]
	\item Según ISO/IEC/IEEE 29148 (6.3.3.3), identificar las necesidades de las partes interesadas incluye: a) solo lo que el cliente dicta; \textbf{b) obtenerlas de las personas y descubrir las necesidades implícitas del dominio y el contexto*}; c) copiar requisitos de sistemas similares; d) priorizar con MoSCoW.
	\item El conocimiento tácito es: a) lo que figura en el reglamento; b) un requisito no funcional; \textbf{c) lo que una persona sabe y usa pero no expresa porque le parece obvio*}; d) un dato confidencial.
	\item La norma 29148 advierte que consultar a una sola clase de partes interesadas: \textbf{a) produce información sesgada*}; b) acelera el proyecto sin riesgos; c) es lo recomendado en proyectos pequeños; d) reemplaza la observación.
	\item ¿Cuál es una pregunta abierta y neutral? a) ¿Verdad que una app sería mejor?; b) ¿El sistema es lento?; \textbf{c) ¿Cómo registra hoy una entrega, desde que llega el socio hasta que se le paga?*}; d) ¿Cuánto tarda el pesaje y quién revisa la humedad?
	\item El cuestionario es más útil para: a) descubrir requisitos desconocidos al inicio; \textbf{b) medir cuántas personas comparten algo ya identificado*}; c) acordar requisitos entre roles; d) observar el trabajo real.
	\item La diferencia principal entre un taller de requisitos y un grupo focal es que el grupo focal: a) dura varios días; b) lo conduce el gerente; \textbf{c) busca conocer percepciones y no tomar decisiones*}; d) no necesita moderador.
	\item La técnica más indicada cuando el trabajo es rutinario y las personas no saben explicarlo es: a) el cuestionario; \textbf{b) la observación*}; c) la lluvia de ideas; d) la lectura de normas.
	\item En la lectura basada en perspectivas: a) un solo lector revisa todo el documento; \textbf{b) se lee el mismo documento desde el punto de vista de distintos usuarios, cada uno con sus preguntas*}; c) se leen solo los títulos; d) se reemplaza a las entrevistas.
	\item ¿Cuál pregunta «¿Cómo podríamos\ldots?» está bien formulada? a) ¿Cómo podríamos hacer una app móvil para el acopio?; \textbf{b) ¿Cómo podríamos lograr que el técnico registre una entrega sin hacer esperar al socio más de cinco minutos?*}; c) ¿Cómo podríamos usar inteligencia artificial?; d) ¿Cómo podríamos terminar el proyecto?
	\item En el repositorio público del PFC se registra: a) el audio de cada entrevista; b) los consentimientos firmados; \textbf{c) un manifiesto con archivo, fecha, duración y huella SHA-256*}; d) los nombres y teléfonos de los participantes.
\end{enumerate}

% ============================================================
\seccion{Observaciones de alineación con el sílabo}
% ============================================================

\begin{enumerate}[leftmargin=20pt,itemsep=4pt]
	\item \textbf{Material anterior no alineado.} Las diapositivas \texttt{ISR401\_\allowbreak{}Semana04.pptx} (carpetas \texttt{ISR401\_\allowbreak{}Diapositivas}, \texttt{\_\allowbreak{}v2} y \texttt{Semanas01-18\_\allowbreak{}Mejoradas}) y la lectura \texttt{Semana\_\allowbreak{}04\_\allowbreak{}ISR401.pdf} tratan prototipado, i*, KAOS, BPMN e historias de usuario y anuncian una «PE-4»; corresponden al plan anterior y no a la semana 4 del sílabo. Este paquete las reemplaza. Lo mismo ocurre con la semana 3 anterior (técnicas clásicas de elicitación), que en el sílabo es contexto y perspectivas.
	\item \textbf{Sumativas vinculadas a GA-TL-02.} La guía dice «Prepara: GA-PE-02, TA-PFC-02», conforme a las notas de alineación del 4 de noviembre; el cronograma del encuadre v3 dice «(TA-PFC-E1, GA-PE-02)». Como GA-TL-02 y TA-PFC-E1 cierran el mismo sábado, la formativa no puede ser requisito de E1; conviene corregir el encuadre a «GA-PE-02, TA-PFC-02».
	\item \textbf{Publicación de TA-PFC-E1.} La guía indica publicación el lunes 9 de noviembre; la tabla de tareas del encuadre indica el lunes 16. La publicación anticipada favorece a los estudiantes; basta unificar la fecha en el encuadre.
	\item \textbf{Referencia IREB del proyecto.} La entrada \texttt{ireb\_\allowbreak{}cpre\_\allowbreak{}fl\_\allowbreak{}2026} del \texttt{referencias.bib} de las guías lista solo a Bühne y Glinz; el temario 3.3.0 (1 de abril de 2026) tiene cuatro autores: Bühne, Glinz, van Loenhoud y Staal. La entrada está corregida en \texttt{referencias\_\allowbreak{}semana04.bib} y conviene reemplazarla en el \texttt{referencias.bib} de las guías. La entrada de ISO/IEC/IEEE 29148 se completó con su DOI.
	\item \textbf{Horario.} Los tiempos asumen una sesión de 3 horas y otra de 2; si el paralelo tiene otra distribución, mover el taller completo a la sesión práctico-experimental y conservar su orden.
\end{enumerate}

\nocite{iso29148_2018,ireb_cpre_fl_2026,bano_2019,pohl_2025}
\bibliographystyle{unsrtnat}
\bibliography{referencias_semana04}

\end{document}
